
%\needspace{4\baselineskip}
\color{uniblue}{\subsection{Полуавтоматическое включение (ПАВ)}}
\color{black}

\begin{enumerate}[label=\arabic{section}.\arabic{subsection}.\arabic{enumi}, labelsep=4pt, leftmargin=0pt, itemindent=57pt, itemsep=0pt, parsep=0pt] % если что можно крутить itemsep

\item
В устройстве предусмотрен режим полуавтоматического включения (ПАВ). Данная функция предполагает включение выключателя с контролем синхронизма или улавливанием синхронизма при условии оперативного включения выключателя на противоположном конце линии. Для ввода режима ПАВ В1 предусмотрен входной сигнал ФСУ \textit{«ПАВ В1»}. При оперативном включении линии с противоположного конца и выполнении условий синхронизма происходит включение выключателя. Если в течение времени \textbf{«Тожид ПАВ В1»} не будут выполнены условия синхронизма, то схема ПАВ В1 выведется из работы.

\item
Функциональная схема логики ПАВ В1 приведена на рисунке \ref{pict_pav:PAV}. Функциональная схема логики ПАВ В2 выполняется аналогично схеме логики ПАВ В1.

\vspace{2mm}
\begin{figure}[H]
\centering
\includegraphics[width=1\textwidth,height=1\textheight,keepaspectratio]{\fbpath/761.3061 ПАВ/img/PAV.pdf}
\captionof{figure}{Функциональная схема логики ПАВ В1}\label{pict_pav:PAV}
\end{figure}

\end{enumerate}